\documentclass{report}

\usepackage[utf8]{inputenc}
\usepackage{textgreek}
\usepackage{blindtext}
\usepackage[inference]{semantic}
\usepackage{multicol}
\usepackage{hyperref}
\usepackage[newfloat]{minted}
\usepackage{amssymb}
\usepackage{amsmath}
\usepackage{graphicx}
\usepackage{svg}
\usepackage[capitalise]{cleveref}
\usepackage{enumitem}
\usepackage{float}
\usepackage{microtype}
\usepackage[a4paper, margin=1in]{geometry}
\usepackage{fancyvrb, newverbs, xcolor}
\usepackage{caption}
\usepackage[UKenglish]{babel}
\usepackage[parfill]{parskip}

\definecolor{cverbbg}{gray}{0.93}
\newcommand{\vtilde}{$\sim\ $}

\newenvironment{equ}{\vspace{0mm}\captionsetup{type=listing}}{\vspace{2mm}}
\SetupFloatingEnvironment{listing}{name=Code Snippet}

\newcommand{\spipe}{\quad\lvert\quad}
\newcommand{\ctexttt}[1]{\colorbox{cverbbg}{\texttt{#1}}}
\newverbcommand{\cverb}
  {\setbox\verbbox\hbox\bgroup}
  {\egroup\colorbox{cverbbg}{\box\verbbox}}

\hypersetup{
    colorlinks = true,
    urlcolor = blue,
    linkcolor = blue,
    citecolor = red
}

\newcommand\MyBox[1]{
  \fbox{\lower0.75cm
    \vbox to 1.7cm{\vfil
      \hbox to 1.7cm{\hfil\parbox{1.4cm}{#1}\hfil}
      \vfil}%
  }%
}

\newcommand{\hs}{\mintinline{haskell}}
\newcommand{\mono}{\mintinline{text}}

\begin{document}

%Start of titlepage!
\begin{titlepage}
    \centering
    
    \textls{\textsc{\huge{Utrecht University}}}\\
    \vspace{5em}
    \textsc{\Large{Thesis Report}}\\
    \vspace{2em}
    \textsc{\large{Department of Information and Computing Sciences}}\\
    
    \vspace{3em}
    \hrulefill
    \vspace{4em}
    
    \huge{Embracing Core \& Supervising the Optimisation Pipeline}\\
    \vspace{1em}
    \Large{Empowering Haskell developers with a looking glass into the core2core pipeline}

    \vspace{4em}
    \hrulefill
    \vspace{3em}
   
    \begin{tabular}{p{0.5\textwidth}r}
        \textit{Author:} & \textit{Supervisors:}\\
        H.A. Peters, BSc & W. Swierstra\\
        h.a.peters@uu.nl & T. McDonall\\
        5917727 &
    \end{tabular}
    
    \vspace{3em}
    \textsc{December 2022}\\
    % \textsc{\textcolor{red}{Version: Draft 2}}
    
    \begin{figure}[b]
        \hspace{-1.25cm}
        \includegraphics{figs/UU_logo_2021_NL_RGB.png}
    \end{figure}
    
\end{titlepage}

\vspace*{\fill}
\begin{center}
\begin{minipage}{.7\textwidth}
\centerline{\textbf{Abstract}}
Straight away communicate the idea that core exploration should be easier and more approachable and give
and drill it home with a small but daunting core snippet.
 
Haskell is a language explicit about side effects which allows for pretty drastic semantic preserving transformations,
inviting programmers to hand over a lot of optimisation responsibilities. To reasure success one needs to explore
interactions with the compiler.


\end{minipage}
\end{center}
\vfill % equivalent to \vspace{\fill}

\thispagestyle{empty}

\newpage
\clearpage
\pagenumbering{roman}
\tableofcontents

\newpage
\clearpage
\pagenumbering{arabic}

%Add with newpage
\newpage

\chapter{Introduction}

\section{GHC, an optimising compiler}

Haskell is a high level language designed to write highly composable functions.
This seemingly encourages programmers to write code that
is not necessarily fast to evaluate. An extremely common example is the composition
of list operations. Consider the function \mono{halves} which divides each element
in a list of \mono{Int}s in halve, discarding those that are not a multiple of 2.

\begin{listing}[H]
\begin{minted}[linenos]{haskell}
halves :: [Int] -> [Int]
halves = map (`div` 2) . filter even
\end{minted}
\end{listing}

Both \mono{map} and \mono{filter} are defined as a loop over an input list, producing a new list as output. 
If code was generated for \mono{halves} in its current form, it would produce two loops as well as allocate
an intermediate list. This is unnecessary extra work and allocation behavior because
it would be possible to rewrite the function to circumvent this issue:

\begin{listing}[H]
\begin{minted}[linenos]{haskell}
halves_fast :: [Int] -> [Int]
halves_fast [] = []
halves_fast (x:xs) = 
  let 
    tl = halves_fast xs 
  in if even x 
     then (x `div` 2):tl
     else tl
\end{minted}
\end{listing}

But requiring the programmer to do such rewrites manually tragically undermines
the benefits of this compositional style of programming. Code simply would be harder to read, write,
maintain, and consequently tend to be less correct. 

Luckily, GHC does address this issue - and many others - with an extensive set of optimisation transformations.
This particular program will benefit greatly from the fusion system which specifically deals with removing
the intermediate lists. This is a well-established optimisation that is also referred to as deforestation \cite{WADLER1990231}.
As a result, compiling with these optimisations enabled will result in a \textbf{syntactically} equivalent
definition for \mono{halves} and \mono{halves_fast}! 

But this poses a question of trust; No compiler is every perfect, so how can we be sure that our code is
correctly optimised and will continue to be in the future? Furthermore, if optimisations are missed, how
can we diagnose the root of the problem and explain what went wrong?

\section{The cascade effect}

Optimisation transformations are applied is a certain order, giving rise to the \textit{cascade effect} \cite{haskell_optimisations_1997}.
This effect refers to the dramatic consequences that the order of transformations can have. We will study an example that
showcases a problematic tug-of-war between two optimisations: (1) inlining and (2) rewrite rules.

\subsection{The inlining transformation}

Inlining is arguably the most important optimisation for any functional language. Coincidentally,
it is also one of the simplest transformations to implement and comprehend. To reveal
why it is such a staple, we must consider it in conjunction with \textbeta-reduction. 
These transformations are fairly simply defined as:

\begin{listing}[H]
\begin{minted}[linenos]{haskell}
-- inlining, knowing that x = x' from some binding site
e             -inline->   e[x := x']

-- beta reduction:
(\x -> e) a   -ß->        e[x := a]
\end{minted}
\end{listing}

We can apply this in a Haskell context to get a more familiar perspective. 
Consider the boolean negation function \mono{not} and a basic usage thereof:

\begin{listing}[H]
\begin{minted}[linenos]{haskell}
not :: Bool -> Bool
not = \x -> if x then False else True

top = True
bottom = not top
\end{minted}
\end{listing}

In the case of \hs{bottom} function we can elect to inline both \mono{not} and \mono{top}
in its body. This gives us 

\begin{listing}[H]
\begin{minted}[]{haskell}
bottom = (\x -> if x then False else True) True
\end{minted}
\end{listing}

We can now perform \textbeta-reduction by taking the body of the lambda function
and substituting \mono{x} by \mono{True}:

\begin{listing}[H]
\begin{minted}[]{haskell}
bottom = if True then False else True
\end{minted}
\end{listing}

Finally, we eliminate the \mono{if} statement by evaluation. Knowing that it always takes the first branch, we
get the final result of:

\begin{listing}[H]
\begin{minted}[]{haskell}
bottom = True
\end{minted}
\end{listing}

We can now see that thanks to this transformation duo, we have reduced an expression to a constant
expression, eliminating any runtime cost. Because the inlining transformation so commonly goes hand
in hand with \textbeta-reduction we will --- for brevity's sake --- from now on presume that \textbeta-reduction,
wherever relevant, is also applied when we say that an inlining transformation has taken place.

In a sense we were lucky here that the body of the \mono{not} function reduced so completely. But in reality
it can be the case that function definitions are quite large. While inlining could then still be a good idea (as it eliminates jumps),
it does come at the cost of larger code sizes. GHC uses quite a few crude heuristics to help with weighing this trade-off \cite{haskell_optimisations_1997}.
Interestingly, these heuristics have not undergone much reconsideration until quite recently \cite{inling_magic_numbers}. 
Suffice to say that it is not trivial to predict whether a function will be inlined and the performance implications thereof.

\subsubsection{Laziness and thunks}
Being a lazy language, Haskell can profit from inlining in a secondary way as well. A well-established
language feature is the let expression, a syntactically nice way bind extra definitions. Consider a function
that decrypts the password of a user but only if it is the admin.

\begin{listing}[H]
\begin{minted}[linenos]{haskell}
getAdminPassword :: User -> Maybe String
getAdminPassword user = 
  let decrypted :: String
      decrypted = decryptSHA1 (password user)

  in case (username user) of
    "admin" -> Just decrypted
    _       -> Nothing
\end{minted}
\end{listing}

Obviously the \mono{decryptSHA1} is going to be an extremely cost heavy function that under no
circumstances we wish to be evaluated unless absolutely necessary. Luckily, because Haskell is a
lazy language the let assignment does not actually evaluate anything; it only allocates a so-called
\textit{thunk}. Such a thunk represents an expression that is yet to be evaluated. Evaluation is
in fact deferred to the point where the thunk is actually needed (which is never if the username is
not \hs{"admin"}). So, because of this property \mono{getAdminPassword} is actually quite effective 
at avoiding unnecessary work. 

Yet it can still be improved upon. While the let expression might not do any extra work it does require
this thunk allocation, which is not free. But because our let bound variable \mono{decrypted} 
is used exactly once, we can inline it at the only call-site without running the risk of duplicating 
the expensive operation. This yields the more performant definition:

\begin{listing}[H]
\begin{minted}[linenos]{haskell}
getAdminPassword :: User -> Maybe String
getAdminPassword user = case (username user) of
    "admin" -> Just (decryptSHA1 (password user))
    _       -> Nothing
\end{minted}
\end{listing}

\subsection{Rewrite rules}
\label{section:introduction:rewrite_rules}

We've seen how generic program transformation may improve performance. However, GHC can only use relatively
shallow reasoning as to not jeopardize the correctness of transformations or explode compile times. 
The programmer on the other hand may have much more in-depth knowledge about the domain of the program 
and its intended behavior. \cite{playing_by_the_rules}.
Programmers can leverage this knowledge by defining so-called \textit{rewrite rules}. They inform the compiler of
substitutions that are not obvious or strictly speaking even correct.  

Consider the binary tree datatype \mono{Tree} with along with a the higher order \mono{mapTree} functions that
facilities transforming the values contained in the \mono{Leaf}s. We then use \mono{mapTree} to compose two
traversals with an addition in the function \mono{addFive}. Obviously that function is rather contrived as
an example of something that someone would write, but the pattern may very well show up during the transformation pipeline.


\begin{listing}[H]
\begin{minted}[linenos]{haskell}
data Tree a = Leaf a | Node (Tree a) (Tree a) deriving Show

mapTree :: (a -> b) -> Tree a -> Tree b
mapTree f (Leaf x) = Leaf (f x)
mapTree f (Node lhs rhs) = Node (mapTree f lhs) (mapTree f rhs)

addFive :: Tree Int -> Tree Int
addFive = mapTree (+1) . mapTree (+4)
\end{minted}
\end{listing}

By now it should be clear why \mono{addFive} is non-optimal; it has a superfluous traversal and allocates an
intermediate structure. \hs{mapTree (+5)} is the far superior equivalent. However, for GHC to infer this fact
it has to perform a too complicated and generic analysis. But we can add the rewrite rule \textit{mapTree/mapTree}
to convince GHC that consecutive applications of \mono{mapTree} are allowed to fuse:

\begin{minted}[linenos]{haskell}
{-# Rules
   "mapTree/mapTree"   forall f g.   mapTree f . mapTree g    = mapTree (f. g)   ;
#-}
\end{minted}

But as it turns out, \textit{mapTree/mapTree} never fired for a reason that is not immediately obvious.

\subsubsection{A note on common abstraction}

A clever reader might have noticed that \mono{mapTree} is a perfect candidate for an implementation of \mono{fmap}
as part of the \mono{Functor} typeclass: 

\begin{minted}[linenos]{haskell}
instance Functor Tree where
  fmap = mapTree
\end{minted}



\subsection{Tug-of-war}
\label{section:introduction:tug_of_war}

A common manifestation of the cascade effect is the tug-of-war between inlining and the application
of rewrite rules. Continuing with the \mono{Tree} example from the previous section, we find that losing this
tug-of-war is the exact reason that the \textit{mapTree/mapTree} rule never fired. Because the inlining operation
was performed first, the left-hand-side of \textit{mapTree/mapTree} no longer occurs in the program and the rule
is rendered non-active. The final optimised function has regrettably converged to the following form:

\begin{minted}[linenos]{haskell}
addFive :: Tree Int -> Tree Int
addFive = \x -> mapTree (+1) (mapTree (+4) x)
\end{minted}

The important point is here is not that rewrite rules are inherently flawed (after all, we could easily
drum up a secondary rule that does fire here), but that one optimisation
may open or close the door to many other optimisations down the road.
From this cascade effect follows that the interaction of a Haskell program
with the optimiser may be quite unstable and consequently sensitive to small changes. 
Changes not only in the source but also in the build environment (a minor release of
the compiler comes to mind). Thus, we cannot trust that our successfully optimised program will remain
optimised in the future. We observe that each program we write may require specific, manual, 
effort to be made more efficient.

\subsection{Non-functional requirements: \mono{inspection-testing}}
\label{section:introduction_inspection_testing}

So if we want our programs to not only be correct but also terminate in reasonable amount of time while
not consuming an overly large chunk of resources, we have to identify an extra set of contraints. 
These constraints do not deal with the functionality of the program but rather it's compiled form.
This collection of additional constraints are examples of \textit{Non-functional} requirements. 

To illustrate with a real world example, the very popular \mono{text} library makes the following promise:
`\textit{Most of the functions in this module are subject to fusion, meaning that a pipeline of such functions will usually allocate at most one Text value.}'
\cite{inspection_testing}.
Like with \mono{addFive} in \cref{section:introduction:tug_of_war}, 
such promises cannot be checked with traditional tests as they do not concern the functionality of the code. 

And as identified by Breitner \cite{inspection_testing}, the aforementioned promise by the \mono{text} library had in fact
been broken in version \mono{1.2.3.2}, shown by the following counter example:

\begin{minted}[linenos]{haskell}
import qualified Data.Text as T
import qualified Data.Text.Encoding as TE
import Data.ByteString

countChars :: ByteString -> Int
countChars = T.length . T.toUpper . TE.decodeUtf8
\end{minted}

Although \mono{countChars} uses a value of type \mono{Text} during the computations, it does not
need to be actually constructed in the final composite definition. As we learn from the definition:

\begin{minted}{haskell}
data Text = Text ByteString Int Int
\end{minted}

\mono{Text} text is merely a \textit{view} into a \mono{ByteString}
by virtue of an offset and length parameter. This means that length reducing operations can cleverly avoid the
costly task of modifying the underlying \mono{ByteString} and instead just change the offset and length parameters.
Now because UTF-16 contains surrogate pairs, the character length of a \mono{Text} value cannot directly be determined
from the byte-length its \mono{ByteString} and still requires $O(n)$ time. However, this does not justify constructing
a concrete \mono{Text} value as opposed to just using the \mono{ByteString} and the offset and length parameters as
separate bindings.

An analogous situation would be a factory with production line workers that
pack and unpack their intermediate results between every exchange in the assembly line. While they may
conveniently communicate about receiving a `car door, a bolt, and a nut' unified as their `input', they
never actually mean to suggest that they wish to receive them packaged together. So too with our \mono{Text} values,
it is mighty handy to communicate about a package of a \mono{ByteString} and two \mono{Int}s, but we never intent to
actually package them at every turn.

But as mentioned, despite the extensive set of rewrite rules that the \mono{text} library has, the
ideal compilation result was not achieved. In itself this example formed the main motivation to 
develop a method to tests these non-functional requirements.
The result is the \mono{inspection-testing} package \cite{inspection_testing}. 
It provides the machinery necessary to add the following statement
directly to the source file, preventing the same regression from occurring in the future.

\begin{minted}{haskell}
inspect $ 'countChars `hasNoType` ''T.Text
\end{minted}

Despite preemptively saving us from future regressions, \mono{inspection-testing} does not help to
identity and path underlying cause. Consider when the test fails at
some point, and you are tasked with finding the root of the problem. In this very example the final 
optimised definition of \mono{countChar} will have undergone many expanding transformations producing 
over 100 lines to painfully sift through without little more ergonomics than a string search. 

It would be possible to get the output at various intermediate moments to gauge where the 
offending \mono{Text} should have been erased, resulting in an extra compilation artifact per transformation.
Usually this means you have to analyze around 20 different versions of the code. 
This is a tedious and error-prone process, not to mention requires a relatively highly skilled programmer 
with a solid understanding of the GHC and its optimiser. 

Moreover, we risk having to repeat this work in the future when any small number of changes
could --- through the cascade effect --- trigger this test as failing again.

\section{Introducing \mono{HsComprehension}}

We intend to address the tediousness and skill required for exploring interactions of specific 
pieces of Haskell source code with GHC's optimisation pipeline.

We believe that there is an opportunity to improve the way that Haskell programmers reason about
the performance characteristics of their code while simultaneously appealing to a larger audience
of less experienced programmers.

This belief stems first from the results of the yearly
Haskell survey where in 2021 only 16\% of respondents either agreed or strongly agreed with the
statement `\textit{I can easily reason about the performance of my Haskell programs}' \cite{haskell_survey_2021}.
We are not the only one seeing this statistic as a call to action. As recently the current year,
Young announced work on a complete Haskell optimisation handbook \cite{optimisation_handbook}.
A preliminary version of the book already shows that a section on analysing optimisation transformations
is planned. Although, it is our opinion that a guiding resource --- while certainly extremely helpful ---
is not a substitute for better tooling.

Secondly, through conversations with Haskell programmers at Well-Typed --- who are also major contributors
and maintainers of GHC itself -- we learned that it was common practice to create personal tools that
assisted them in analysing intermediate results of the optimisation transformations during the fairly
frequent task of trying to make performance critical sections of programs more efficient. 
This proves that there is a demand for such tools and that a unified solution does not yet exist.



\chapter{Research Questions}


\paragraph{Main Question} 
How can GHC’s core2core passes be captured and presented in such a way that core
inspection becomes a more successful, and easier task?

\hfill \break

\paragraph{Sub-Question 1} 
How does one efficiently navigate where small changes occur in two or more captures?

\paragraph{Sub-Question 2}
How to make viewing core more manageable using various display options?

\paragraph{Sub-Question 3}
How could performance regressions that have occurred in the past in popular Haskell projects,
have been resolved faster?

\chapter{Background}

\section{A Core Language}
\label{section:background:core_lang}

GHC recognises three distinct stages of the compilation process \cite{haskell_optimisations_1997}.

\begin{enumerate}
  \item \textbf{Frontend}, parsing and type checking, enables to give clear errors that reference the unaltered source code.
  \item \textbf{Middle}, a number of optimisation transformations.
  \item \textbf{Backend}, generates and compiles to machine code, either via C-- or LLVM.
\end{enumerate}

The middle section is tasked with doing all the optimisation work it can, leaving only those optimisations to the backend
it can not otherwise perform. Within the middle section, the work is further split into substages, where each transformation
is a separate, composable pass. An obvious benefit to this approach is that each pass can be tested independently. Moreover, because
the types are preserved throughout the middle section, it can be verified that each transformation preserves type correctness;
Providing some evidence towards the correctness of the transformation.

Regardless, the Haskell source language itself is not a good target for optimisation. The source has to be
rich and expressive, which yields an AST datatype of over 100 constructors. Any transformation over such a datatype has far too
many cases to be considered practical to implement; Not to mention the myriad of bug-inducing edge cases. To overcome this issue, the
middle section first performs a desugaring pass, translating the source language into a far simpler intermediate language
called \textit{Core}. 

Core --- which is how we will refer to GHC intermediate representation going forward --- is designed to be as simple as possible
without being impractical. A testament to that property is the fact that we can fit the entire definition on the page:

\begin{listing}[H]
\begin{minted}{haskell}
data Expr
  | Var Id
  | Lit Literal
  | App Expr Expr
  | Lam Bndr Expr
  | Let Bind Expr
  | Case Expr Bndr Type [Alt]
  | Cast Expr CoercionR -- safe coercions are required for GADTs
  | Tick Tickish Expr  -- annotation that can contain some meta data
  | Type Type
  | Coercion Coercion
  
data Alt = Alt AltCon [Bndr] Expr
 
data Bind
  = NonRec Bndr Expr
  = Rec [(Bndr, Expr)]
  
-- the Var type contains information about
-- a variable, like it's name, a unique identifier
-- and analysis results
type Bndr = Id
type Id = Var

type Program = [Bind]
\end{minted}
\label{code:core_def}
\caption{An ever so slightly simplified version of the Core Language}
\end{listing}

Writing an optimisation transformation essentially of type \hs{Program -> Program} does not now seem
as daunting, given the drastically reduced number of cases to consider. 
Likewise, maintaining and debugging transformations is much less of a strenuous task.

\section{The core2core transformations}

Over its numerous decades of development, the core2core pipeline has been fitted with a myriad of transformations.
It would be impractical to discuss all of them here. However, we will discuss in depth the role of the multiple simplifier
passes, as well as the worker/wrapper transformation. The first because it gives context to the rewrite rules and the latter
because it is a natural bridge to unboxed types; both concepts which will become relevant in discussing the results of this thesis.
Lastly, we zoom in on the analysis results stored in the \mono{Var} type.

\subsection{The simplifier: its parts}

The simplifier is quite literally an indispensable part of the transformation pipeline. Although the parts of the pipeline
are meant to be separable entities, they heavily rely on the simplifier to deal with the cleanup. As such, it has earned
itself the reputation of being the workhorse of the pipeline \cite{ghc_wiki_core2core}.

If you were looking for an atom, you have not found it. The simplifier is in itself again a collection of smaller transformations,
albeit applied repeatedly and interleaved such that they cannot be easily untangled. Each simplifier subpart is a local transformation,
that is, it only looks at the immediate surroundings of the expression. We give a near comprehensive list of each subpart along with an example: 
\cite{ghc_wiki_core2core}

\subsubsection{Constant Folding}
Evaluate expressions that require no runtime information. This is a very logical transformation that does not
require any further considerations.

\begin{listing}[H]
\begin{minted}{haskell}
-- before
3 + 4
-- after
7
\end{minted}
\end{listing}

\subsubsection{Inlining}
\label{section:background:inlining}

Inlining function calls is a well known performance enhancing operation in all language, but especially so in functional ones.
The difference is staggering with a around 10-15\% performance increase for conventional languages, as opposed to
30-40\% for functional languages \cite{haskell_optimisations_1997}.

Unlike constant folding, it is not an ad-hoc no-brainer. An easy case would be to inline a function that is
used exactly once. All that this does is remove the veil that hides the function body for further optimisations:

\begin{listing}[H]
\begin{minted}{haskell}
-- before
f x = x + 1
f 3
-- after
3 + 1
\end{minted}
\end{listing}

However, if the function is used in multiple locations with different arguments, you risk increasing the size of the program
because the function body will be duplicated at each call-site. This is a trade-off that --- although often worth it --- must be
considered on a case to case basis. GHC has a number of heuristics to determine whether inlining is worth it or not. \cite{secrets_of_the_inliner}
Obviously, this cannot be a perfect solution and one can imagine how a small change in the code can suddenly cause the inliner to reverse
its decision; a testament to the non-continuous nature of the compiler.

Besides inlining function calls, Haskell's let bindings also form an opportunity for inlining.
After all, let bindings are often described as syntax sugar for lambda abstractions, but there is an important difference to be considered.
Because let bindings in Haskell are lazily evaluated, it may lead to explosion if the let bound variable is used in a lambda
expression. For example:

\begin{listing}[H]
\begin{minted}{haskell}
-- before, 'expensive' is called at most once
let v = expensive 42
    l = \x -> ... v ...
in map l xs

-- after, 'expensive' is called for each element of 'xs'
let l = \x -> ... expensive 42 ...
in map l xs
\end{minted}
\end{listing}

In this case, inlining would be disastrous for performance and GHC will take great care to avoid it. 

\subsubsection{Case of known constructor}

The case destruction is only the way to get to the WHNF of an expression. This means that inside of a case expression
we have learned information about the variable under scrutiny and can use it to infer the result of outer case expressions.
Consider the following scenario:

\begin{listing}[H]
\begin{minted}{haskell}
safe_tail :: [a] -> [a]
safe_tail ls = case ls of
    [] -> []
    x:xs -> tail ls
\end{minted}
\end{listing}

Which inlining will transform into:

\begin{listing}[H]
\begin{minted}{haskell}
safe_tail :: [a] -> [a]
safe_tail ls = case ls of
    [] -> []
    x:xs -> case ls of
        [] -> error "tail of empty list"
        (x:xs) -> xs
\end{minted}
\end{listing}

Since we scrutinize \mono{ls} again in the inner case, we actually already know that the second case is only ever called,
thus we can eliminate the case expression as a whole.

\begin{listing}[H]
\begin{minted}{haskell}
safe_tail :: [a] -> [a]
safe_tail ls = case ls of
    []   -> []
    x:xs -> xs
\end{minted}
\end{listing}

The removal of the call \mono{error} during this process is a testament to this function being deserving of the \mono{safe} prefix.

\subsubsection{Case of case}

There are cases (no pun intended) where the case-of-know-constructor cannot quite do its job, although it is obvious
that benefits are still to be gained. Consider for example what happens when instead of scrutinizing the same variable,
the outer case scrutinizes the result of the innner case:


\begin{listing}[H]
\begin{minted}{haskell}
-- before (possibly desugared from 'if (not x) then E1 else E2'
-- after also inlining 'not')
case (case x of {True -> False; False -> True}) of
    True -> E1
    False -> E2
\end{minted}
\end{listing}

We might hope to gain something from the information that the inner case has produced by moving the outer case expression
to each branch of the inner one:

\begin{listing}[H]
\begin{minted}{haskell}
case x of
    True -> case False of {True -> E1; False -> E2}
    False -> case True of {True -> E1; False -> E2}
\end{minted}
\end{listing}

Now we can rely on constant folding to simplify all the way down to:

\begin{listing}[H]
\begin{minted}{haskell}
case x of
    True -> E2
    False -> E1
\end{minted}
\end{listing}

Note that we do risk duplicating \mono{E1} and \mono{E2} if the expression does not reduce so completely.
A solution to this are \textit{join points} \cite{compiling_wo_continuations},
which we will not go into here for the sake of brevity.

\subsubsection{Rewrite rules}

We already discussed in \cref{section:introduction:rewrite_rules} how rewrite rules are way to express domain
specific knowledge to the compiler that it otherwise can not practically infer. Applying given rewrite rules is task
also given to the simplifier. It should now be clear that this process can get a little messy since the simplifier
is under no obligation to apply the rules, nor its other tasks, in any particular order. We will get into this issue
more in the next section where we discuss the simplifier at large.

It should be noted that the use of rewrite rules are very common during the compilation of most any Haskell program.
This is because the internal fusion system on list operations are implemented as built in rewrite rules.
We discuss this system more in depth in \cref{section:background:fusion}.



\subsection{The simplifier: its sum}

An attentive reader might have noticed that when explaining one part of the simplifier, we often relied on another to
its job. This begs the question: how does the simplifier determine in which order to run its subparts. The answer to that is
that it does not. The analogy here is that a compiler is a gun and a program is a target. Every program is very different
and we cannot know in advance how to best hit it. Therefore, we must ensure that the compiler --- or this case the simplifier stage ---
has a lot of bullets in its gun to ensure being able to effectively hit a lot of targets \cite{haskell_optimisations_1997}.

A secondary problem with this approach is that one simplifier part may enable another to its job after. Running the simplifier
once is therefore not satisfactory, and we must run it until it reaches a sort of \textit{fixed point}; At least up until some
arbitrary limit since there is no guarantee such a fixed point exists and even if it does that we find it and not get stuck in a loop.

Aforementioned looping behavior can actually occur quite easily due to certain rewrite rules. It is not always the case that
the RHS of a rewrite rules objectively better than the right. It might be that the rewrite is benificial because it \textbf{may}
enable other optimisation to take place afterwards. However, if for some reason that did not turn out to be possible, we may want
to apply the reverse rewrite rule later. This is obviously problematic as we can be ping-pong between rewrite rules forever.
To overcome that issue one can enable/disable rewrites rules in certain phases of simplifier. To understand this we must first
understand when the simplifier is run.

In order, the simplifier is --- rather unintuitively --- interspersed between other transformations as follows:

\begin{enumerate}
  \item \textbf{Gentle} (disables case-of-case transformations)
  \item \textbf{Phase 2}
  \item \textbf{Phase 1}
  \item \textbf{Phase 0}
  \item \textbf{Final}
\end{enumerate}

By default, rewrite rules as well as inlinings can occur in each of these phases but the programmer does have the ability
to specify deviations from this behavior. For example, in the \mono{text} library, we find in the \mono{Data.Text} module
the following snippet:

\begin{listing}[H]
\begin{minted}{haskell}
-- This function gives the same answer as comparing against the result
-- of 'length', but can short circuit if the count of characters is
-- greater than the number, and hence be more efficient.
compareLength :: Text -> Int -> Ordering
compareLength t c = S.compareLengthI (stream t) c
{-# INLINE [1] compareLength #-}

{-# RULES
"TEXT compareN/length -> compareLength" [~1] forall t n.
    compare (length t) n = compareLength t n
#-}
\end{minted}
\end{listing}

Here phase control is used to indicate that \mono{compareLength} should \textbf{only} be inlined in phase 1 and conversely
that the rewrite rule \textbf{compareN/length} may be applied \textbf{except} in phase 1. What this ensures is that the inliner
will not operate on the result of the rewrite rule directly in the same phase. This is supposedly because we expect \mono{compareLength}
to occur in the LHS of a different rewrite rule which is to be desired over inlining at first.

\subsection{The worker/wrapper transformation}
\label{section:background:worker_wrapper}

The \textit{worker/wrapper} transformation is able to change the type of computation (typically a loop) into a \textit{worker} function
along with a \textit{wrapper} that can convert back to the original type. The example listed on the GHC wiki is that of the reverse
function on lists \cite{ghc_wiki_wwrapper}. One might concisely define with direct recursion and the \mono{++} operator:

\begin{listing}[H]
\begin{minted}[linenos]{haskell}
reverse :: [a] -> [a]
reverse [] = []
reverse (x:xs) = reverse xs ++ [x]
\end{minted}
\end{listing}

Of course this is not very efficient as the \mono{++} operator itself is already linear, make the reverse function quadratic; If we
create an auxiliary \textbf{worker} function \mono{reverse'} that takes an extra accumulator argument, we can create a linear version:

\begin{listing}[H]
\begin{minted}[linenos]{haskell}
reverse :: [a] -> [a]
reverse = reverse' []
  where reverse' acc [] = acc
        reverse' acc (x:xs) = reverse' (x:acc) xs
\end{minted}
\end{listing}

Here the \textbf{wrapper} is simply applying the empty list as a starting argument to the function. Thus, by introducing some state
that exists only during the lifetime of the computation, the function has become asymptotically more efficient.

Impressively, the transformation may also greatly improve the constant factor of the runtime by leveraging unboxed types.
Unlike strict languages, the \mono{Int} type in Haskell --- while always representing a 64 bit integer --- is not statically sized.
After all, any value is lazy and may therefore still be unevaluated thunk whose size is unknown at compile time. As a result, \mono{Int}s
are always stored on the heap and thus require runtime allocation. However, \mono{Int} has a strict counterpart \mono{Int#} (unboxed integer)
in which it is defined: \hs{data Int = I# Int#}. Computations on \mono{Int#}s can be evaluated much more cheaply since such values can be
stored on the stack.

Knowing this, we can opt to create a worker that does add state to the computation, but changes the types to their unboxed counterparts.
Naturally the wrapper operation is then simply the constructor of the lazy type \hs{I#}. Let us consider the example of the recursive \mono{triangular} function which
given a number \mono{n} returns the sum of all numbers from 1 to \mono{n}:

\begin{listing}[H]
\begin{minted}[linenos]{haskell}
triangular :: Int -> Int
triangular 0 = 0
triangular n = n + triangular (n-1)
\end{minted}
\end{listing}

Although an all-knowing compiler could have determined that the result of \mono{triangular} can simply be calculated numerically in constant time as $\frac{n^2 + n}{2}$,
we can still trust GHC to infer an important property about the function using strictness analysis. Namely: \mono{triangular} does not produce any intermediate results that can
used without evaluating the whole of \mono{triangular}. That is: if you are willing to spent any amount of time on \mono{triangular}, you might as well calculate the whole thing.
Thus GHC decides to rewrite the function using a strict worker, removing the need for many dubious allocations:

\begin{listing}[H]
\begin{minted}[linenos]{haskell}
wtriangular :: Int# -> Int#
wtriangular 0 = 0
wtriangular n = GHC.Prim.+# n (wtriangular (GHC.Prim.-# n 1))

triangular :: Int -> Int
triangular (I# w) = I# (wtriangular w)
\end{minted}
\end{listing}

To get a feeling for the difference in performance, we can compare the runtime of the original function with the worker/wrapper version:

\begin{table}[H]
\begin{tabular}{|l|l|}
\textbf{Snippet} (compiled with -O0) & \textbf{Result of \mono{time} with input \mono{10000000}}\\
\hline \\
\mono{original} & 0,10s user 0,03s system 99\% cpu 0,132 total \\
\mono{transformed} & 0,01s user 0,01s system 98\% cpu 0,023 total \\
\end{tabular}
\caption{The runtime of both version of \mono{triangular}. We can see that the worker/wrapper tranformation has increased runtime performance by a significant factor.}
\end{table}

It might not be immediately intuitive why the performance differs so drastically and where exactly these seemingly evil allocation occur.
The ridiculousness of the original function becomes apparent when we consider a C implementation with the same allocation behavior. 
Although lacking in laziness, we can consider an \mono{Int} to map to a \mono{long*} (pointing to heap allocated memory) 
and an \mono{Int#} to map to a plain \mono{long}.


\begin{minted}[linenos]{c}
#include <stdio.h>
#include <stdlib.h>

// utility function for heap allocating a 64 bit integer
long* alloc_long() { return (long*)malloc(sizeof(long)); }

long* triangular(const long* n_ptr) {
  // allocate the return value
  long* ret = alloc_long();
  // derefence the pointer into a value
  long n = *n_ptr;
  if (n==0) { *ret = 0; } else {
    // allocate a new pointer with which to call the function recursively
    long* inp_ptr = alloc_long();
    *inp_ptr = n-1 ;
    *ret = n + *triangular(inp_ptr);
  }
  return ret;
}
\end{minted}

Rest assured that any C programmer suggesting this implementation would get some very weird looks.

\subsection{Analysis transformations}

Consider again the Core ADT given in \cref{section:background:core_lang}, it was mentioned that the \mono{Var}
type is used to represent variables and their various properties. We look into the one field of \mono{Var} that
is dynamic during the core2core pipeline: \mono{IdInfo}.

This data type follows the concept of weakening of information, i.e. the information is never incorrect but
may be less precise or even missing. Furthermore, the \mono{IdInfo} may differ for different \mono{Var} instances
that refer to the same variable.

It is as of \mono{ghc-9.4.2} defined as:

\begin{listing}[H]
\begin{minted}[linenos,fontsize=\footnotesize]{haskell}
data IdInfo
  = IdInfo {
        ruleInfo        :: RuleInfo,
        -- ^ Specialisations of the 'Id's function which exist.
        -- See Note [Specialisations and RULES in IdInfo]
        realUnfoldingInfo   :: Unfolding,
        -- ^ The 'Id's unfolding
        inlinePragInfo  :: InlinePragma,
        -- ^ Any inline pragma attached to the 'Id'
        occInfo         :: OccInfo,
        -- ^ How the 'Id' occurs in the program
        dmdSigInfo      :: DmdSig,
        -- ^ A strictness signature. Digests how a function uses its arguments
        -- if applied to at least 'arityInfo' arguments.
        cprSigInfo      :: CprSig,
        -- ^ Information on whether the function will ultimately return a
        -- freshly allocated constructor.
        demandInfo      :: Demand,
        -- ^ ID demand information
        bitfield        :: {-# UNPACK #-} !BitField,
        -- ^ Bitfield packs CafInfo, OneShotInfo, arity info, LevityInfo, and
        -- call arity info in one 64-bit word. Packing these fields reduces size
        -- of `IdInfo` from 12 words to 7 words and reduces residency by almost
        -- 4% in some programs. See #17497 and associated MR.
        --
        -- See documentation of the getters for what these packed fields mean.
        lfInfo          :: !(Maybe LambdaFormInfo),

        -- See documentation of the getters for what these packed fields mean.
        tagSig          :: !(Maybe TagSig)
    }
\end{minted}
\end{listing}

Through non-structure-altering transformations like \textit{Demamd Analysis} and \textit{Strictness Analysis} the \mono{IdInfo}
record is updated accordingly. This information may then be used by future transformations that can only optimise safely or
productively under certain circumstances (remember the disastrous work duplication in \cref{section:background:inlining}?).

The information contained in \mono{IdInfo} is a major source of complexity when it comes to comprehending Core. Consider
for example this `pretty'-printed \mono{IdInfo} instance:

\begin{listing}[H]
\begin{minted}[linenos,fontsize=\footnotesize]{haskell}
[GblId,
 Arity=4,
 Str=<L,U(U(U(U(C(C1(U)),A,C(C1(U)),C(U),A,A,C(U)),A,A),1*U(A,C(C1(U)),A,A),A,A,A,A,A),U(A,A,C(U),...etc.
 Unf=Unf{Src=<vanilla>, TopLvl=True, Value=True, ConLike=True,
         WorkFree=True, Expandable=True, Guidance=IF_ARGS [50 0 0 0] 632 0}]
\end{minted}
\end{listing}

We can quickly learn of few things about the variable it describes. First,  it is apparently a function that has
an arity of 4 (i.e. it takes 4 arguments). Secondly, we obtain some of the magic heuristic values involved with inlining
(also known \textit{unfolding} hence \mono{Unf}). However, if we want to diagnose why for example this variable was evaluated
lazily even though it is always used exactly once, we would have to decode the strictness signature under \mono{Str}.
Currently, the best resource for decoding this would be to read the comments in the GHC source code itself.

\section{The fusion system}
\label{section:background:fusion}

Fusion --- the process of combining multiple operations over a structure into a single operation --- is in GHC
implemented using the \textit{build/foldr} idiom. This system was developed as an improvement on deforestation 
(Waddler \cite{WADLER1990231}) which has shortcomings when it comes to recursive functions \cite{shortcut_fusion}.

We borrow an example from Seo \cite{shortcut_fusion_blog} to illustrate how the build/foldr system works. Consider
the example --- coincidentally very similar to \mono{triangular} (\cref{section:background:worker_wrapper}) --- of
summing a list of numbers:

\begin{listing}[H]
\begin{minted}[linenos]{haskell}
foldr (+1) 0 [1..10]
\end{minted}
\end{listing}

Very concise indeed but its execution would be much slower than say, an imperative for loop, because we first create
a list and then consume it right away. If we look inside the definition of \mono{[1..10]} we find the function \mono{from}:

\begin{listing}[H]
\begin{minted}[linenos]{haskell}
from :: (Ord a, Num a) => a -> a -> [a]
from a b = if a > b
           then []
           else a : from (a + 1) b
\end{minted}
\end{listing}

Which is itself a specialisation for lists. We can write a more generic catamorphism that takes any duo of constructors
of type \mono{a -> b -> b} and \mono{b} respectively (previously \mono{:} and \mono{[]}):

\begin{listing}[H]
\begin{minted}[linenos]{haskell}
from' :: (Ord a, Num a) => a -> a -> (a -> b -> b) -> b -> b
from' a b = \c n -> if a > b
                    then n
                    else c a (from' (a + 1) b c n)
\end{minted}
\end{listing}

The glue between \mono{from} and \mono{from'} is the \textit{build} function, which requires the \mono{Rank2Types} language extension
and re-specialises these generic functions back to lists.

\begin{listing}[H]
\begin{minted}[linenos]{haskell}
build :: forall a. (forall b. (a -> b -> b) -> b -> b) -> [a]
build g = g (:) []

from a b = build (from' a b )
\end{minted}
\end{listing}

Thus far it seems we have only introduced more complexity without any apparent benefit. 
However, this is the point that we run into the key idea: \mono{build} is the antithesis of \mono{foldr} such that the following holds:

\begin{listing}[H]
\begin{minted}[linenos]{haskell}
foldr k z (build g) = g k z
\end{minted}
\end{listing}

And this also gives us the main reductive rewrite rule to make this system work.
So to summarize, we have shown that generalizing lists to list producing functions to catamorphisms
over list like arguments, produces a common interface exposing the elimination of code. More concretely,
because we keep \mono{g} as a generic function, we can choose to give it the arguments of the consequent
\mono{foldr} directly and thus eliminate the intermediate list. This is again a situation like the assembly line
workers analogy from \cref{section:introduction:inspection_testing}; It is very productive to communicate about
functions taking and producing lists, but actually doing so in a chained context is like wrapping and unwrapping
boxes at every step.

Analogies aside, observe how our original expression can now completely be reduced to a single integer value
by rewriting \mono{from} with this common interface.

\begin{listing}[H]
\begin{minted}[linenos]{haskell}
foldr (+) 0 (from 1 10)
-- { inline from }
foldr (+) 0 (build (from' 1 10))
-- { apply rewrite rule }
from' 1 10 (+) 0
-- { inline from' }
\c n -> (if 1 > 10
            then n
            else c 1 (from' 2 10 c n)) (+) 0
-- { beta reduce }
if 1 > 10
   then 0
   else 1 + (from' 2 10 (+) 0)
-- { repeat until base case }
1 + 2 + ... + 9 + 10 + 0
-- { constant fold }
55
\end{minted}
\end{listing}

But what if there is no \mono{foldr} in the expression? What if we use simpler functions like \mono{map} or \mono{filter}?
Well, as most introductory Haskell courses are likely to tell you, most list functions can be defined in terms of \mono{foldr}.
Thus, we can similarly use rewrite rules to map all these common list functions to a combination \mono{build} and \mono{foldr}.
This idea also relieves us from the burden of having to define rewrite rules for all the million possible 
combination of lists operations:

\begin{listing}[H]
\begin{minted}[linenos]{haskell}
map f xs    = build (\ c n -> foldr (\ a b -> c (f a) b) n xs)
filter f xs = build (\ c n -> foldr (\ a b -> if f a then c a b else b) n xs)
xs ++ ys    = build (\ c n -> foldr c (foldr c n ys) xs)
concat xs   = build (\ c n -> foldr (\ x y -> foldr c y x) n xs)
repeat x    = build (\ c n -> let r = c x r in r)
zip xs ys   = build (\ c n -> 
    let zip' (x:xs) (y:ys) = c (x,y) (zip' xs ys)
        zip' _ _ = n
    in  zip' xs ys)
[]          = build (\ c n -> n)
x:xs        = build (\ c n -> c x (foldr c n xs))
\end{minted}
\end{listing}

We can see some these functions at work when looking at the definition of \mono{unlines} and its subsequent
reductions through the build/foldr system:

\begin{listing}[H]
\begin{minted}[linenos]{haskell}
unlines :: [String] -> String
unlines ls = concat (map (\l -> l ++ ['\n']) ls)
-- { rewrite concat and map }
unlines ls = build
  (\c0 n0 -> foldr (\xs b -> foldr c0 b xs) n0 ( build
    (\c1 n1 -> foldr (\l t -> c1 (build
      (\c2 n2 -> foldr c2 ( foldr c2 n2 ( build
        (\c3 n3 -> c3 '\n' n3))) l )) t) n1 ls )))
-- { apply rewrite rule foldr/build }
unlines ls = build
  (\c0 n0 ->
    (\c1 n1 -> foldr (\l t -> c1 (build
      (\c2 n2 -> foldr c2 (
        (\c3 n3 -> c3 '\n' n3) c2 n2 ) l)) t) n1 ls)
          (\xs b -> foldr c0 b xs) n0)
-- { beta reduce }
unlines ls = build
  (\ c0 n0 -> foldr (\l t -> foldr c0 t( build
    (\c2 n2 -> foldr c2 (c2 '\n' n2) l))) n0 ls)
-- { apply rewrite rule foldr/build }
unlines ls = build
  (\c0 n0 -> foldr (\l b -> foldr c0 (c0 '\n' b) l) n ls)
-- { inline build }
unlines ls = foldr (\l b -> foldr (:) ('\n' : b) l) [ ] ls
-- { inline foldr }
unlines ls = h ls
    where h [] = []
          h (l:ls) = g l
              where g [] = '\n' : h ls
                    g (x:xs) = x : g xs
\end{minted}
\end{listing}

What we end up with is the most efficient implementation of \mono{unlines} that we could possibly write by hand \cite{shortcut_fusion}.
Because the list \text{ls} is the input of the function we cannot remove it all together, but keep in mind that when inlining call to this function
at the use site, foldr/build fusion may apply again. 

\section{Contemporary Haskell comprehension}

\subsection{Communicating in Core}
\label{section:communicating_core}

We are not pioneers discovering the land of Core inspection. Since its inception, Core has
been used to communicate about programs and compiler interactions. Sifting through open issues on the
GHC compiler itself, one quickly comes across discussions elaborated by Core snippets. Consider issue
\href{https://gitlab.haskell.org/ghc/ghc/-/issues/22207}{\#22207} titled `\textit{bytestring Builder performance regressions after 9.2.3}' for example.
\hfill \break

\textit{While testing the performance characteristics of a bytestring patch meant to mitigate withForeignPtr-related performance regressions,
it was noticed that several of our Builder-related benchmarks have also regressed seriously for unrelated reasons.
The worst offender is byteStringHex, which on my machine runs about 10 times slower and allocates 21 times as much when using ghc-9.2.4 or ghc-9.4.2
as it did when using ghc-9.2.3. Here's a small program that can demonstrate this slowdown:}
\hfill \break

The author then provides two snippets containing the final Core representation of \mono{byteStringHex} as produced by the two different GHC version.
Each of these documents contain around 400 lines of unhighlighted code annotated with all available information. And while having all available information
sounds like a good thing (and it is in a way) it poses a serious practicality issue.
Namely: it is exceedingly difficult for a human to read and comprehend a certain aspect of the AST while having to filter out another.
Not to mention, it solidifies reading Core as an activity reserved for only the most expert Haskell developers by scaring others away with
a steep barrier to entry.

\subsection{Current tooling}

\subsubsection{GHC itself}
\label{section:background:ghc}


Core snippets of your program can easily be coerced out of GHC. The most information you can get
is the Core AST at each pass of the optimisation pipeline by using \mono{-ddump-core2core}.
To reduce the signal-to-noise ratio of Core snippets, one can use any number of suppression options.
It is common to suppress type arguments and type applications for example. These are largely
uninteresting because Core is only typed to facilitate the correctness preserving checks that \mono{core-lint} performs
and do not affect the compilation in any way.

As can be read in the GHC documentation, the following suppression flags are available to help to tame the beast.

\begin{table}[H]
  \begin{tabular}{p{0.35\linewidth}|p{0.65\linewidth}}
  \textbf{GHC flag} & \textbf{Effect on Core printing} \\
  \hline
           & \\
  \mono{-dsuppress-all} & In dumps, suppress everything (except for uniques) that is suppressible.  \\
  \mono{-dsuppress-coercions} & Suppress the printing of coercions in Core dumps to make them shorter  \\
  \mono{-dsuppress-core-sizes} & Suppress the printing of core size stats per binding (since 9.4)  \\
  \mono{-dsuppress-idinfo} & Suppress extended information about identifiers where they are bound  \\
  \mono{-dsuppress-module-prefixes} & Suppress the printing of module qualification prefixes  \\
  \mono{-dsuppress-ticks} & Suppress "ticks" in the pretty-printer output.  \\
  \mono{-dsuppress-timestamps} & Suppress timestamps in dumps  \\
  \mono{-dsuppress-type-applications} & Suppress type applications  \\
  \mono{-dsuppress-type-signatures} & Suppress type signatures  \\
  \mono{-dsuppress-unfoldings} & Suppress the printing of the stable unfolding of a variable at its binding site  \\
  \mono{-dsuppress-uniques} & Suppress the printing of uniques in debug output (easier to use diff)  \\
  \mono{-dsuppress-var-kinds} & Suppress the printing of variable kinds\\
\end{tabular}
\end{table}

We can show how these suppression options greatly improve the readability of Core snippets by comparing the
desugared (unoptimized) Core of \mono{quicksort} with and without the \mono{-dsuppress-all} flags.

First the source:

\begin{listing}[H]
\begin{minted}[linenos]{haskell}
quicksort :: Ord a => [a] -> [a]
quicksort [] = []
quicksort (x:xs) = quicksort (filter (< x) xs) ++ [x] ++ quicksort (filter (>= x) xs)
\end{minted}
\label{code:quicksort_src}
\end{listing}

The desugared Core without suppression:

\begin{listing}[H]
\begin{minted}[linenos,fontsize=\footnotesize]{text}
-- RHS size: {terms: 55, types: 47, coercions: 0, joins: 0/10}
quicksort :: forall a. Ord a => [a] -> [a]
[LclIdX]
quicksort
  = \ (@ a_a1pH) ($dOrd_a1pJ :: Ord a_a1pH) ->
      let {
        $dOrd_a1w8 :: Ord a_a1pH
        [LclId]
        $dOrd_a1w8 = $dOrd_a1pJ } in
      let {
        $dOrd_a1w5 :: Ord a_a1pH
        [LclId]
        $dOrd_a1w5 = $dOrd_a1pJ } in
      \ (ds_d1wq :: [a_a1pH]) ->
        case ds_d1wq of wild_00 {
          [] -> ghc-prim-0.6.1:GHC.Types.[] @ a_a1pH;
          : x_a1hP xs_a1hQ ->
            letrec {
              greater_a1hS :: [a_a1pH]
              [LclId]
              greater_a1hS
                = let {
                    $dOrd_a1pQ :: Ord a_a1pH
                    [LclId]
                    $dOrd_a1pQ = $dOrd_a1pJ } in
                  letrec {
                    greater_a1pT :: [a_a1pH]
                    [LclId]
                    greater_a1pT
                      = filter
                          @ a_a1pH
                          (let {
                             ds_d1wF :: a_a1pH
                             [LclId]
                             ds_d1wF = x_a1hP } in
                           \ (ds_d1wE :: a_a1pH) -> > @ a_a1pH $dOrd_a1pQ ds_d1wE ds_d1wF)
                          xs_a1hQ; } in
                  greater_a1pT; } in
            letrec {
              lesser_a1hR :: [a_a1pH]
              [LclId]
              lesser_a1hR
                = let {
                    $dOrd_a1vV :: Ord a_a1pH
                    [LclId]
                    $dOrd_a1vV = $dOrd_a1pJ } in
                  letrec {
                    lesser_a1vY :: [a_a1pH]
                    [LclId]
                    lesser_a1vY
                      = filter
                          @ a_a1pH
                          (let {
                             ds_d1wD :: a_a1pH
                             [LclId]
                             ds_d1wD = x_a1hP } in
                           \ (ds_d1wC :: a_a1pH) -> < @ a_a1pH $dOrd_a1vV ds_d1wC ds_d1wD)
                          xs_a1hQ; } in
                  lesser_a1vY; } in
            ++
              @ a_a1pH
              (quicksort @ a_a1pH $dOrd_a1w5 lesser_a1hR)
              (++
                 @ a_a1pH
                 (GHC.Base.build
                    @ a_a1pH
                    (\ (@ a_d1wx)
                       (c_d1wy :: a_a1pH -> a_d1wx -> a_d1wx)
                       (n_d1wz :: a_d1wx) ->
                       c_d1wy x_a1hP n_d1wz))
                 (quicksort @ a_a1pH $dOrd_a1w8 greater_a1hS))
        }
\end{minted}
\label{code:quicksort_core_raw}
\end{listing}

The same desugared Core representation with the \mono{-dsuppress-all} flag:

\begin{listing}[H]
\begin{minted}[linenos,fontsize=\footnotesize]{text}
-- RHS size: {terms: 55, types: 47, coercions: 0, joins: 0/10}
quicksort
  = \ @ a_a1pH $dOrd_a1pJ ->
      let { $dOrd_a1w8 = $dOrd_a1pJ } in
      let { $dOrd_a1w5 = $dOrd_a1pJ } in
      \ ds_d1wq ->
        case ds_d1wq of wild_00 {
          [] -> [];
          : x_a1hP xs_a1hQ ->
            letrec {
              greater_a1hS
                = let { $dOrd_a1pQ = $dOrd_a1pJ } in
                  letrec {
                    greater_a1pT
                      = filter
                          (let { ds_d1wF = x_a1hP } in
                           \ ds_d1wE -> > $dOrd_a1pQ ds_d1wE ds_d1wF)
                          xs_a1hQ; } in
                  greater_a1pT; } in
            letrec {
              lesser_a1hR
                = let { $dOrd_a1vV = $dOrd_a1pJ } in
                  letrec {
                    lesser_a1vY
                      = filter
                          (let { ds_d1wD = x_a1hP } in
                           \ ds_d1wC -> < $dOrd_a1vV ds_d1wC ds_d1wD)
                          xs_a1hQ; } in
                  lesser_a1vY; } in
            ++
              (quicksort $dOrd_a1w5 lesser_a1hR)
              (++
                 (build (\ @ a_d1wx c_d1wy n_d1wz -> c_d1wy x_a1hP n_d1wz))
                 (quicksort $dOrd_a1w8 greater_a1hS))
        }

\end{minted}
\label{code:quicksort_core_dsuppress}
\end{listing}

A drastic improvement for sure, but there are still some things to be left desired.
A simpler language like Core generally needs more code to express the same thing, we can thus expect
to generate more data than the original Haskell code. Moreover, should you be interested the state off
the program at each of the intermediate steps, you can expect to see about 20x more data still.
Unless you know exactly what to search for, this begs for a more ergonomic, filtered view of the data.

\subsubsection{GHC Plugins}

GHC --- being also a playground for academics and enthusiasts alike --- is extremely flexible when it comes to
altering its functionality. Using the now well established plugin interface, one is able to hook into almost
any operation of the front- and midsection of the compiler. One such place is managing the core2core passes that
will be run on the current module. This point of customization can be used to intersperse each core2core pass with an
identity transformation that smuggles away a copy of the AST in its full form.

One such existing plugin is \mono{ghc-dump} \cite{ghc_dump}. Besides extracting intermediate ASTs, it defines an
auxiliary Core definition to which it provides a GHC version agnostic mapping. This has the increased benefit of 
being able to directly compare snapshots from different GHC versions; A not uncommon task as discussed in \cref{section:communicating_core}.
And while certainly being an improvement over plain text representation, we believe exploring and comparing such
dumps requires a more rich interface.

\newpage


\chapter{Methods}

\section{Requirements and architecture}

\begin{itemize}
  \item GHC version agnostic (to where reasonable i.e. 8.4 onwards)
  \item Simple and non invase steps to create dumps
  \item Cross-platform ability to explore dumps
  \item Not everything needs to be supported (think unfoldings) but should be extendible in the future
\end{itemize}

We envisioned a high degree of interactability with the snapshots of the intermediate ASTs. To realise this
in a cross-platform fashion, we decided to use a browser based frontend application. Because the concept of
mutually recursive algebraic datatypes are very pervasive in the Core AST, we felt it would be helpful
if the frontend language had first class support for this. This naturally led to us to Elm, a functional language
that compiles to Javscript \cite{elm_lang}.

\section{Creating the GHC plugin}

By far the most logical and convenient method for obtaining all the available information in the AST it to capture
it directly using a GHC core2core plugin that intersperses each transformation with an identity transformation that
saves a JSON representation of the AST to a file as a side effect. 

\subsection{Capturing the information}

If we wish to support multiple recent versions of GHC we need to deal with the fact that the Core ADT has undergone
a few minor changes and additions. We believe that the solution is to create some auxiliary definition to which we can
map various versions of the Core ADT. This was done very efficiently by building upon the existing \mono{ghc-dump} package,
which already implemented such an ADT as well as a version agnostic conversion module with the help of \mono{min_version_GHC}
macro statements \cite{ghc_dump}. 

What \mono{ghc-dump} also intelligently addresses, is the issue of possible infinite recursion
in the Core AST through the \mono{CoreExpr} inside the unfolding of a \mono{CoreBndr}. If we demote each
call site to a simple identifier, we obtain a finite representation that we can later reconstruct by traversing
the AST with an environment.

\subsection{Globally unique variables}

It is not strictly necessary for variable names in a program to be unique. A variable name always
references the nearest binding site. This holds true for Core as well, but is not very convenient
when we want to analyze a certain variable in isolation. After all, naked variables are ambiguous without a given
environment. Consider the definition of tail:

\begin{listing}[H]
\begin{minted}[linenos]{haskell}
tail xs = case xs of
  x:xs -> xs
  _    -> error "tail of empty list"
\end{minted}
\end{listing}

We cannot simply refer to the variable \mono{xs} as that name has two different binding sites.
We solve this by running a \textit{uniqify} pass before each snapshot that freshens all duplicate uniques in the entire
module after each core2core transformation. This gives us the ability to refer to a binding site and its usages unambiguously using simply an integer.
The big idea here is that any viewing logic is completely decoupled from binding semantics:

\begin{listing}[H]
\begin{minted}[linenos]{haskell}
tail xs_0 = case xs_0 of
  x_1:xs_2 -> xs_2
  _        -> error "tail of empty list"
\end{minted}
\end{listing}

It is possible to omit the numbered suffixes when displaying the AST, but internally it is very useful to be able to make
this distinction without any further effort.

\subsection{Detecting changes}

If a module is of a slightly larger size, it becomes difficult to spot the changes made by a certain
transformation, if there even are any. To address this, we decided to develop a feature that allows for
the filtering of code that remains unchanged. Let us define what unchanged means in this context. It is
important to make the subtle distinction between syntactic equivalence and $\alpha$-equivalence. The difference
is that the latter is agnostic to the names of variables, as long as they refer to the same binding site.

We can quickly solve the decision problem of syntactic equality by calculating a hash of an expression beforehand
and simply checking for equality of this hash value. We considered using recent improvements of full sub expression
matching \cite{hashing_mod_alpha}, but decided against it as it was not clear how to effectively present the results
nor did it rarely prove useful to isolate changes in the AST as they were rarely local to begin with.
Instead, we opted for a far simpler approach where we only hash the toplevel functions, and provide a more crude
option to hide any toplevel definitions that have not changed at all.

All in all, we still recommend that issues are attempted to be reproduced in small modules as the amount of
noise can quickly become overwhelming despite change detection.

\section{Creating the frontend application}

\subsection{Reproducing the AST}
It would have been extremely tedious to have to constantly maintain a Core ADT in Elm along with a JSON
parser that is compatible with the JSON output of the Haskell plugin. Luckily, we were able to use the
\mono{haskell-to-elm} \cite{haskell_to_elm} package to automatically derive all the needed boilerplate code.

For example the \mono{Alt} datatype (i.e. an arm of a case expression) is defined as follows in Haskell:

\begin{listing}[H]
\begin{minted}[linenos]{haskell}
data Alt = Alt
    { altCon :: AltCon
    , altBinders :: [Binder]
    , altRHS :: Expr
    }
    deriving (Generic)
\end{minted}
\end{listing}

The corresponding, auto generated, Elm definition becomes:

\begin{listing}[H]
\begin{minted}[linenos]{elm}
type alias Alt =
    { altCon : AltCon
    , altBinders : List Binder
    , altRHS : Expr
    }

altDecoder : Json.Decode.Decoder Alt
altDecoder =
    Json.Decode.succeed Alt |>
    Json.Decode.Pipeline.required "altCon" altConDecoder |>
    Json.Decode.Pipeline.required "altBinders" (Json.Decode.list binderDecoder) |>
    Json.Decode.Pipeline.required "altRHS" exprDecoder
\end{minted}
\end{listing}

Conveniently, it would be very easy in the future to extend the Haskell ADT with additional information
because it is the single source of truth; The Elm ADT and JSON machinery can then be regenerated with a single command.

Additionally, we use this pipeline embellish the AST. Specifically, we add a field to each \mono{BinderId} of the type:

\begin{listing}[H]
\begin{minted}[linenos]{haskell}
type BinderThunk = Found Binder | NotFound | Untouched
\end{minted}
\end{listing}

Intially this field is set to the \mono{Untouched} variant. Once the finitely encoded AST is decoded into the 
Elm universe, a reconstruction traversal takes place that with the help of an environment strengthens the 
\mono{BinderThunk} with a reference to its binder. The \mono{NotFound} variant is pure for verification purposes
and should never occur in a sound input program. It is also important to note that Elm's single static assignment
semantics naturally implies that the \mono{Binder} is never copied, but only referenced.


\subsection{Pretty printing}

It certainly should be considered useful to display the Core in exactly the same way that
GHC does. After all, this is what programmers are currently already used to and its design
has been given a lot of thought. However, we felt a need to first create a separate representation
that is tailored to those who have not seen Core before.

Haskell programmers that care enough to inspect the interaction with the compiler are likely to be
avid readers of at least basic Haskell syntax. Therefore, we decided to create a pretty printer that attempts
to be as similar to Haskell source as possible. Just like GHC, we used pretty printing the method developed 
by Waddler \cite{prettier_printer}, implemented in Elm using \mono{elm-pretty-printer} \cite{prettier_printer_elm}.
We believe that such a representation is a suitable way to minimize the shock reaction that newcomers might have
when they first encounter Core.

To facilitate syntax hightlighting, the pretty printer adds the appropriate token identifiers such that 
\mono{pygments} \cite{pygments} can be used to colorize the output.

To compare, consider again the Core representation presented in \cref{section:background:ghc} with the quicksort
representation produced by our pretty printer:

\begin{listing}[H]
\begin{minted}[linenos]{haskell}
quicksort :: forall a. Ord a -> [a] -> [a]
quicksort a $dOrd ds = case ds of
  { : x xs -> GHC.Base.++ @a
                (quicksort @a $dOrd
                   (GHC.List.filter @a
                      (\ds -> GHC.Classes.< @a $dOrd ds x) xs))
                (GHC.Base.++ @a
                   (GHC.Base.build @a
                      (\a c n -> c x n))
                   (quicksort @a $dOrd
                      (GHC.List.filter @a
                         (\ds -> GHC.Classes.> @a $dOrd ds x) xs)))
    [] -> GHC.Types.[] @a
  }
\end{minted}
\end{listing}

Notice that although this might look like normal Haskell, it contains explicit type variables like
\mono{a}, hinting at the SystemFC nature of Core.

\subsection{Including the source}
It goes without saying that the source code of a module is an essential part of any analysis. Therefore,
the plugin copies the source code and runs it through the \mono{pygmentize} \cite{pygments} tool to obtain
an HTML representation of the source code that is highlighted exactly the same way as the pretty printed Core.
This HTML source code is embedded in the output of the dump.

The artificacts produced by a single module now look like this:

\begin{listing}[H]
\begin{minted}[linenos]{text}
Quicksort_0.json
Quicksort_10.json
Quicksort_11.json
Quicksort_12.json
Quicksort_13.json
Quicksort_14.json
Quicksort_15.json
Quicksort_16.json
Quicksort_17.json
Quicksort_18.json
Quicksort_1.json
Quicksort_2.json
Quicksort_3.json
Quicksort_4.json
Quicksort_5.json
Quicksort_6.json
Quicksort_7.json
Quicksort_8.json
Quicksort_9.json
Quicksort.hs
Quicksort.html
Quicksort_meta.json
\end{minted}
\end{listing}

All these files will be compressed in a zip archive to keep the size of the output smaller. It also makes for
a convenient atomic entity to load up in the frontend application where it is unzipped on the fly.

\subsection{Unfoldings and capture sizes}

The entire body of a variable is sometimes saved as part of its \mono{IdInfo} which is exported to interface
files such that inlining can take place across modules. This part of the \mono{IdInfo} struct is called the
\textit{unfolding} of a variable. By nature, this can be a very large expression. At this time, we have no
need for the exact unfolding in the frontend application, other than knowing it exists. Therefore, we decided 
to replace any unfolding with a string literal to indicate that the unfolding was removed:

\begin{listing}[H]
\begin{minted}[linenos]{haskell}
removeUnfolding :: Unfolding -> Unfolding
removeUnfolding u@CoreUnfolding {..} = 
  u { unfTemplate = ELit (MachStr (T.pack "unfolding removed by plugin")) }
removeUnfolding x = x
\end{minted}
\end{listing}

This reduces the size of captures in some cases by a factor of 2.

\subsection{Improving performance with cache semantics}

The way that Elm operates --- like many other interactive browser based applications --- is by rendering the HTML on every
update and then diffing the result with the current state of the DOM to create a minimal change list. This is generally a very
good strategy because mutating the DOM is by far the most dominant cost contributor. However, when it comes to pretty printing a
large Core expression the \mono{view} function starts to incur a significant cost. This is wasteful because many updates do not
actually affect the output of the pretty printer (events like \mono{onHover} and \mono{onClick} for example).
Initially this led to some serious usability issues where the application would stutter and freeze.

We were able to overcome this problem using the \mono{Html.Lazy} module, which takes some HTML producing closure and its arguments
separately. If on some update the arguments have not changed, the closure is not evaluated and a cached result is returned. Regardless,
the aforementioned diffing still takes place, reducing the cost of the update to negligible. Of course, we cannot get around the fact that
any updates that do affect the output of the pretty printer might still cause some stutters, but the frequency of such updates is 
generally far lower.

\subsection{Deployment}

By virtue of being solely a stateless HTML + JavaScript application after compilation, means that the frontend can
easily and cheaply be deployed to any static file hosting service. Because the dump files are never send to the server,
we can discard any privacy concerns while still providing a no effort method to analyze the dumps. Anyone is still free
of course, to build and host their own build of the frontend which is similarly open sourced.


\newpage



\chapter{Results}

\section{Tour of the tool}
Basic outline to prepare for the 'experiments' ahead

\section{Lists vs Streams}
Here we discuss the basic of the stream paper, how it operates, why haskell has moved on by using join points and
finally how the compilation pipeline treats these 2.

\section{Diagnosing a failing inspection test}

Hearkening back to \mono{inspection-testing} discussed in \cref{section:introduction_inspection_testing}, we put
ourselves in the shoes of a programmer who gets surprised by a failing inspection test and reason how she might
employ HsComprehension to diagnose the problem. She breaks the process down into a number of steps.

\paragraph{1. Isolate the problem}
Modules typically contain more than 1 function and during the core2core transformations many more auxiliary functions are
added and many functions are inlined to produce exponentially more code. To get good data it is helpful to temporarily isolate the
failing test case into a separate module. She is able to reproduce the following minimal case:


\begin{minted}{haskell}
{-# LANGUAGE TemplateHaskell #-}

module InspectionTests where

import Test.Inspection
import qualified Data.Text as T
import qualified Data.Text.Encoding as TE
import Data.ByteString

countChars :: ByteString -> Int
countChars = T.length . T.toUpper . TE.decodeUtf8

-- the failing test case
inspect $ 'countChars `hasNoType` ''T.Text
\end{minted}

The following error is produced by the build:

\begin{minted}{sh}
app/InspectionTests.hs:21:1: countChars `hasNoType` Data.Text.Internal.Text failed:
# ...
# 700 lines of core as seen at the end of the core2core pipeline
# ...
\end{minted}

So despite isolating the problem, the amount of textual data generate is still unwieldy.
It is also incomplete in the sense that it does not show the process that produced this final
artefact.

\paragraph{2. Creating a dump}


% Trying to reach the same conclusion as https://github.com/haskell/text/issues/202 (recursing into the mentioned https://github.com/haskell/text/pull/348)


\section{Performance regression at Channable}

%https://github.com/haskell/text/issues/202
%https://github.com/haskell/text/pull/348
%
%Another performance regression: https://gitlab.haskell.org/ghc/ghc/-/issues/22207


\chapter{Discussion}

\section{Review of the tech stack}
Elm does not provide amazing performance given the domain (mostly due to being strict),
but it does make it easy to write haskell like code over an AST and thus requires less skill
for other haskell programmers to contribute to. The final result is also just a website which is
very minimal barrier to entry and can be used to share things.

\section{Printing Core like GHC}
If we truly wish to be a useful tool for GHC developers and experts alike, we must respect
that the current core printer form is well established and should provide at least a mode in which
the tool prints in a very similar manner. This does not imply that a different beginner mode, as you
could call the current state of affairs, that prints a more readable, slightly desugared version, 
is not also invaluable in itself.
\section{Post Core factors}
Core2Core isn't everything, like in this issue: https://github.com/channable/alfred-margaret/issues/24.
Things like STG and the backend are also part of the equation.

\bibliographystyle{abbrv}
\bibliography{refs}

\appendix
\clearpage
\pagenumbering{roman}
\chapter{Appendix}
An appendix here!

\end{document}
